\chapter{研究の新規性分析と総合演習}
\label{chap10}

\begin{chapterbox}
\textbf{【この章で学ぶこと】}
\begin{itemize}
\item AIを使って自分の研究の新規性を客観的に分析・言語化できる
\item これまでの章の成果物を統合した研究ポートフォリオを完成させる
\item AI活用の全体像を振り返り、今後の活用方針を立てられる
\end{itemize}
\end{chapterbox}
\vspace{5mm}

\section{なぜ新規性の言語化が重要か}
\index{しんきせい@新規性}

研究の価値は、突き詰めれば「何が新しいのか」に集約される。しかし、多くの大学院生にとって、自分の研究の新規性を明確に言語化することは難しい。新規性の言語化が求められる場面は、研究生活の中で繰り返し登場する。

\subsection*{論文投稿（査読対応）}

学術論文を投稿すると、査読者から必ず問われるのが「この研究の新規性は何か（What is the novelty of this work?）」である。Introductionで先行研究との差分が明確に記述されていなければ、どれほど実験結果が良くても、論文は採択されない。

\subsection*{学会発表（Introductionでの位置づけ）}

学会発表のスライドでは、冒頭のIntroductionで「既存研究の何が足りなくて、自分の研究が何を解決するのか」を明確に示す必要がある。

\subsection*{就職面接での研究説明}

第\ref{chap9}章で述べた通り、就職面接では「あなたの研究のオリジナリティは何ですか？」という質問が頻出である。

\subsection*{研究計画書の作成}

博士課程への進学を志望する場合や、科研費などの研究費申請を行う場合、研究計画書の中で「提案する研究の新規性」を説得力を持って記述しなければならない。

これらすべての場面において、\textbf{自分の研究の新規性を客観的に分析し、適切な言葉で表現する能力}が求められる。

\section{新規性分析の3つの観点}

研究の新規性は、一つの軸だけで語れるものではない。ここでは、新規性を分析するための3つの観点を紹介する。

\subsection{観点1: 方法論的新規性 --- 新しい手法・アプローチの提案}
\index{ほうほうろんてきしんきせい@方法論的新規性}

最も典型的な新規性の形態である。既存の手法では解決できない問題に対して、新しいアルゴリズム、フレームワーク、モデル、評価手法などを提案するものである。

方法論的新規性を主張する際の表現パターン:
\begin{itemize}
\item ``We propose a novel approach that combines X and Y to address Z''
\item ``Unlike existing methods that ..., our method ...''
\item ``To the best of our knowledge, this is the first work that ...''
\end{itemize}

\subsection{観点2: 実践的新規性 --- 実世界への適用・検証}
\index{じっせんてきしんきせい@実践的新規性}

必ずしも手法そのものが新しくなくても、既存手法を新しい領域・対象に適用し、その有効性を検証することは重要な新規性である。

実践的新規性を主張する際の表現パターン:
\begin{itemize}
\item ``We apply X to the domain of Y for the first time''
\item ``We validate the effectiveness of X through a real-world case study of Y''
\item ``Our empirical evaluation reveals that ...''
\end{itemize}

\subsection{観点3: 時間的観点 --- 最新の研究動向の中での位置づけ}

研究は常に時間軸の中に存在する。最新の技術動向、規格の改定、社会的ニーズの変化に対応することも、重要な新規性の一形態である。

時間的観点での新規性を主張する際の表現パターン:
\begin{itemize}
\item ``With the recent advancement of X, there is an urgent need for Y''
\item ``In light of the latest revision of standard Z, we propose ...''
\item ``The emerging challenge of X necessitates a new approach to Y''
\end{itemize}

\subsection*{3つの観点の関係}

多くの優れた研究は、これら3つの観点のうち複数を兼ね備えている。

\begin{table}[H]
\centering
\begin{tabular}{lll}
\toprule
観点 & 問いかけ & 強い場合の特徴 \\
\midrule
方法論的新規性 & 手法のどこが新しいか？ & 新アルゴリズム、新フレームワーク \\
実践的新規性 & どこに適用したのが新しいか？ & 新領域への適用、大規模実験 \\
時間的観点 & なぜ「今」この研究が必要か？ & 最新動向への対応、新規格対応 \\
\bottomrule
\end{tabular}
\end{table}

\section{AIを使った新規性分析ワークフロー}

自分の研究の新規性を客観的に分析するために、3つのステップからなるワークフローを提案する。

\subsection{Step 1: 研究の位置づけマッピング}

\begin{promptbox}
\# 役割\\
あなたは[研究分野名]の専門家であり、学術論文の査読経験が豊富な研究者です。

\# 依頼\\
以下の研究の学術的な位置づけを分析してください。

\# 研究の概要
\begin{itemize}
\item 分野: [例: 自動運転の安全性工学]
\item テーマ: [例: GSNと安全状態レポートを統合した安全性保証手法の提案]
\item 主要なキーワード: [例: Safety Case, GSN, Autonomous Driving, ODD]
\item 使用した手法: [例: GSNによる安全ケース構築、安全状態レポートとの統合]
\item 主な結果: [例: ケーススタディで有効性を検証]
\end{itemize}

\# 分析してほしい内容
\begin{enumerate}
\item この研究が属する学術分野・サブ分野の俯瞰図
\item 関連する研究コミュニティ（学会、ジャーナル）
\item この分野の主要な研究潮流（3〜5つ）
\item 各潮流における自分の研究の位置づけ
\item 隣接する研究分野との接点
\end{enumerate}
\end{promptbox}

\textbf{注意}: AIの回答はあくまで「仮説」として扱うこと。特に、具体的な論文名や著者名が挙げられた場合は、必ず実際にその論文が存在するか確認する必要がある（第3章で学んだハルシネーション対策）。

\subsection{Step 2: 既存研究との差分分析}

\begin{promptbox}
\# 役割\\
あなたは学術論文の査読者です。

\# 依頼\\
以下の「自分の研究」と「先行研究」を比較し、差分を分析してください。

\# 分析してほしい内容
\begin{enumerate}
\item 各先行研究と自分の研究の共通点
\item 各先行研究と自分の研究の相違点
\item 先行研究全体に共通する「限界」で、自分の研究が解決しているもの
\item 比較表（行: 比較項目、列: 各研究）
\end{enumerate}
\end{promptbox}

このような比較表は、論文のIntroductionやRelated Workセクションで非常に効果的である。AIに下書きを作成させた後、自分で正確性を検証し、必要に応じて修正する。

\subsection{Step 3: 新規性の言語化}

\begin{promptbox}
\# 役割\\
あなたは学術英語ライティングの専門家です。

\# 依頼\\
新規性分析の結果を、英語論文のIntroductionの最後に置く「本研究の貢献」パラグラフとしてまとめてください。

\# 条件
\begin{itemize}
\item 3文で構成（各文が1つの貢献に対応）
\item 文1: 方法論的新規性（何を提案したか）
\item 文2: 実践的新規性（何で検証したか）
\item 文3: 時間的観点（なぜ今重要か）
\item ``The main contributions of this paper are as follows:'' で始めること
\end{itemize}
\end{promptbox}

\section{総合演習: 研究ポートフォリオの完成}

本章の後半は、これまでの章で作成してきた成果物を統合し、研究ポートフォリオとして完成させる総合演習である。

\subsection*{これまでの成果物の振り返り}

\begin{table}[H]
\centering
\begin{tabular}{lll}
\toprule
章 & 成果物 & 内容 \\
\midrule
第3章 & 文献調査レポート & 先行研究のサーベイ \\
第4章 & 英語Abstract & 英文要旨（150〜250語） \\
第5章 & 発表スライド & 学会発表やゼミ発表用 \\
第6章 & 研究コード & 実験・分析に使用したコード \\
第7章 & データ分析・グラフ & データの可視化と分析結果 \\
第8章 & GitHub Webページ & 研究成果を公開するWebページ \\
第9章 & ガクチカ & 就職活動用のES文章 \\
第10章 & 新規性分析レポート & 研究の新規性を整理した文書 \\
\bottomrule
\end{tabular}
\end{table}

\subsection*{GitHubリポジトリでの整理と公開}

これらの成果物を、一つのGitHubリポジトリとして体系的に整理する。

\begin{verbatim}
my-research-portfolio/
├── README.md               ← リポジトリの概要
├── docs/
│   ├── literature-review/  ← 文献調査レポート
│   ├── abstract/           ← 英語Abstract
│   ├── novelty-analysis/   ← 新規性分析レポート
│   └── job-search/         ← ガクチカ（公開範囲に注意）
├── slides/                 ← 発表スライド
├── src/                    ← 研究コード
├── data/                   ← データ分析・グラフ
├── .gitignore
└── LICENSE
\end{verbatim}

\section{AI活用の振り返りと今後}

\subsection*{AIが特に効果的だった場面}

\begin{table}[H]
\centering
\begin{tabular}{lll}
\toprule
場面 & 効果 & 理由 \\
\midrule
文章の構造化 & 非常に高い & パターン認識と整理が得意 \\
英語表現の改善 & 高い & 大量の英語テキストで学習 \\
コードの雛形生成 & 高い & 定型的なパターンの再現が得意 \\
アイデアの壁打ち & 高い & 多角的な視点を提供できる \\
想定質問の生成 & 高い & 多様なパターンを網羅的に生成 \\
\bottomrule
\end{tabular}
\end{table}

\subsection*{AIの限界が明らかだった場面}

\begin{table}[H]
\centering
\begin{tabular}{lll}
\toprule
場面 & 限界 & 対処法 \\
\midrule
事実の正確性 & ハルシネーションのリスク & 必ず原典で検証 \\
最新情報の把握 & 学習データの時点で止まる & 最新論文は自分で確認 \\
研究の独創性 & 既存パターンの組み合わせが限界 & 本質的な着想は人間の役割 \\
実験の設計 & 分野固有の暗黙知が不足 & 指導教員や先輩に相談 \\
感情・動機の表現 & 表面的になりがち & 自分の言葉で書く \\
\bottomrule
\end{tabular}
\end{table}

\subsection*{今後のAI進化への対応}

\textbf{ツールは変わっても、原理は変わらない}。本書で学んだ以下の原理は、どのようなAIツールが登場しても適用できる。

\begin{itembox}[l]{AI活用の4つの普遍的原理}
\begin{enumerate}
\item \textbf{明確な指示の重要性} --- ROCF（Role, Output, Context, Format）の考え方は、あらゆるAIツールに応用できる。
\item \textbf{段階的な作業プロセス} --- 複雑なタスクを段階に分けて進めることで、人間が各段階で検証・修正できる。
\item \textbf{批判的検証の習慣} --- 99\%正確なAIの残り1\%の誤りを見つけられるのは、専門知識を持つ人間だけである。
\item \textbf{人間の判断が必要な領域の認識} --- 研究の方向性を決める、仮説を立てる、実験結果を解釈する、倫理的な判断をする——これらはAIが「代替」できない。
\end{enumerate}
\end{itembox}

\subsection*{「AIを使いこなす」研究者になるために}

AIの普及により、「AIを使える人」と「使えない人」の間に生産性の格差が生まれている。しかし、より本質的な格差は、\textbf{「AIに使われる人」と「AIを使いこなす人」}の間にある。

\textbf{「AIに使われる人」の特徴}:
\begin{itemize}
\item AIの出力をそのまま使う
\item AIなしでは作業できなくなる
\item AIの回答が正しいかどうか判断できない
\end{itemize}

\textbf{「AIを使いこなす人」の特徴}:
\begin{itemize}
\item AIの出力を批判的に検証し、取捨選択する
\item AIを使って自分の思考を加速・拡張する
\item AIの限界を理解し、適切な場面で活用する
\end{itemize}

AIは道具である。素晴らしい道具であるが、あくまで道具である。あなたの研究は、あなたにしかできない知的営みである。AIはその営みを支える強力な助手であり、本書がその助手との付き合い方を学ぶ一助となれば幸いである。

\vspace{5mm}
\begin{itembox}[l]{本章のキーポイント}
\begin{enumerate}
\item \textbf{新規性は3つの観点で分析する} --- 方法論的新規性、実践的新規性、時間的観点
\item \textbf{新規性分析は段階的に行う} --- 位置づけマッピング→差分分析→言語化の3ステップ
\item \textbf{成果物は統合して初めて価値が最大化する} --- GitHubリポジトリとして体系的に整理する
\item \textbf{ツールは変わっても原理は変わらない} --- 明確な指示、段階的プロセス、批判的検証、人間の判断
\item \textbf{「AIを使いこなす」研究者を目指す} --- AIは助手であり、研究の本質は人間の知的営みにある
\end{enumerate}
\end{itembox}

\section*{演習問題}

\begin{enumerate}
\item[\textbf{演習10-1}] \textbf{新規性の3観点分析}\\
自分の研究（または取り組んでいるテーマ）について、方法論的新規性、実践的新規性、時間的観点の3つの観点から新規性を分析しなさい。各観点について、3〜5文で記述すること。

\item[\textbf{演習10-2}] \textbf{先行研究との比較表の作成}\\
自分の研究テーマに関する先行研究を3本以上取り上げ、比較表の形式で自分の研究との差分を整理しなさい。比較項目は5つ以上設定すること。AIを使って下書きを作成してもよいが、各セルの内容の正確性は自分で検証すること。

\item[\textbf{演習10-3}] \textbf{貢献記述の作成}\\
演習10-1と10-2の結果をもとに、英語論文のIntroductionで使える「本研究の貢献」パラグラフ（3文）を作成しなさい。日本語訳も併記すること。

\item[\textbf{演習10-4}] \textbf{研究ポートフォリオの構築}\\
本書の各章で作成した成果物を含むGitHubリポジトリを作成しなさい。README.mdを英語で記述し、リポジトリの構成、各成果物の概要、実行方法（コードがある場合）を記載すること。

\item[\textbf{演習10-5}] \textbf{AI活用方針書の作成}\\
本書で学んだ内容を踏まえ、今後の研究活動におけるAI活用方針書（A4用紙1枚程度）を作成しなさい。以下の項目を含めること。
\begin{itemize}
\item 自分の研究活動でAIを活用する場面のリスト
\item 各場面でのAI活用の具体的な方法
\item AIに任せない（人間が判断する）領域の明記
\item AI使用の倫理的ガイドライン（自分なりの基準）
\item 半年後に見直す予定日
\end{itemize}
\end{enumerate}
